\documentclass[11pt,a4paper]{article}\usepackage[utf8]{inputenc}

\usepackage[T1]{fontenc}
\usepackage{tgpagella}
\usepackage[spanish]{babel}
\usepackage{geometry}
\usepackage{graphicx}
\usepackage{booktabs}
\usepackage{array}
\usepackage{xcolor}
\usepackage{titlesec}
\usepackage{fancyhdr}
\usepackage[parfill]{parskip}
\usepackage[expansion=false]{microtype}
\usepackage{hyperref}
\usepackage{setspace}
\usepackage{needspace}

\geometry{top=3.0cm,bottom=3.0cm,left=3.5cm,right=3.5cm}
\setlength{\parskip}{0.65em}
\setlength{\parindent}{0em}

\definecolor{azulai}{RGB}{30,80,140}
\definecolor{doradoai}{RGB}{140,90,0}
\definecolor{grisai}{RGB}{70,70,70}

\titleformat{\section}{\Large\bfseries\color{azulai}}{}{0em}{}[{\color{azulai}\titlerule[0.5pt]\nopagebreak[4]}]
\titlespacing*{\section}{0pt}{1.8em}{0.8em}
\titleformat{\subsection}{\large\bfseries\color{doradoai}}{}{0em}{}[{\nopagebreak[4]}]
\titlespacing*{\subsection}{0pt}{1.4em}{0.5em}
\titleformat{\subsubsection}{\normalsize\bfseries\color{grisai}}{}{0em}{}[{\nopagebreak[4]}]
\titlespacing*{\subsubsection}{0pt}{1.0em}{0.3em}

\pagestyle{fancy}\fancyhf{}
\rhead{\textcolor{grisai}{\small Izas — Arquitectura Interna}}
\lhead{\textcolor{grisai}{\small Carta Natal Completa}}
\cfoot{\textcolor{grisai}{\small\thepage}}
\renewcommand{\headrulewidth}{0.3pt}

\hypersetup{colorlinks=true,linkcolor=azulai,urlcolor=azulai}
\setstretch{1.45}
\tolerance=1500
\emergencystretch=4em

\begin{document}

% ── Portada ──────────────────────────────────────────────────────────────────
\begin{titlepage}
  \centering
  \vspace*{1.5cm}
  {\Huge\bfseries\color{azulai} Carta Natal Completa}\\[0.5cm]
  {\large\color{grisai} Arquitectura Interna}\\[0.3cm]
  {\small\itshape\color{grisai} Un método para sostener cuerpo, energía y vida con coherencia}\\[2cm]
  {\huge\color{doradoai} Izas}\\[1.5cm]
  {\Large 16/04/1979 \quad 04:20}\\[0.3cm]
  {\Large Pamplona, España}\\[0.3cm]
  {\normalsize Lat: 42.8157° \quad Lon: -1.6522° \quad UTC+02:00}\\[0.3cm]
  {\normalsize Ascendente: Acuario 6°20' \quad
    MC: Escorpio 29°12'}\\[2.5cm]
  \vfill
  {\small Generado el 19/05/2026}
\end{titlepage}

\tableofcontents
\newpage

% ── Rueda ─────────────────────────────────────────────────────────────────────
\section{Carta Natal}
\begin{figure}[h!]
  \centering
  \includegraphics[width=0.90\textwidth]{Izas_AI_rueda.png}
  \caption{Carta natal de Izas — 16/04/1979 04:20 — Pamplona, España}
\end{figure}

\newpage

% ── Tabla de posiciones ───────────────────────────────────────────────────────
\section{Posiciones Planetarias}

\begin{center}
\begin{tabular}{llrrlll}
  \toprule
  \textbf{Planeta} & \textbf{Signo} & \textbf{Grado} & \textbf{Casa} &
  \textbf{Elemento} & \textbf{R} \\
  \midrule
  Sol & Aries* & 25°31' & 2 & Fuego &  \\
  Luna & Sagitario & 8°24' & 10 & Fuego &  \\
  Mercurio & Piscis & 28°57' & 2 & Agua &  \\
  Venus & Piscis & 21°30' & 1 & Agua &  \\
  Marte & Aries* & 7°01' & 2 & Fuego &  \\
  Júpiter & Cáncer & 29°42' & 6 & Agua &  \\
  Saturno & Virgo & 7°33' & 7 & Tierra & R \\
  Urano & Escorpio & 19°58' & 9 & Agua & R \\
  Neptuno & Sagitario & 20°21' & 11 & Fuego & R \\
  Plutón & Libra* & 17°38' & 8 & Aire & R \\
  Quirón & Tauro & 9°16' & 3 & Tierra &  \\
  Lilith & Leo & 20°38' & 7 & Fuego &  \\
  Nodo Norte & Virgo & 16°50' & 7 & Tierra &  \\
  Nodo Sur & Piscis & 16°50' & 1 & Agua &  \\
  \bottomrule
\end{tabular}
\end{center}

\vspace{0.5cm}
\begin{itemize}
  \item \textbf{Ascendente:} Acuario 6°20'
  \item \textbf{Medio Cielo:} Escorpio 29°12'
  \item \textbf{Signos interceptados} (marcados con * en la tabla):
  \begin{itemize}
  \item Aries (casa 2) -- Libra (casa 8)
  \end{itemize}
\end{itemize}

\newpage

% ── Interpretación ────────────────────────────────────────────────────────────
\section{Interpretación — Arquitectura Interna}

\begin{center}
{\small\itshape
No se trata de definirte. Se trata de observar cómo se organiza tu energía,\
dónde puedes perder equilibrio y qué apoyos te ayudan a sostenerte.
}
\end{center}

\vspace{0.5cm}

\subsection{1. Visión General}

Tu carta muestra la siguiente distribución de elementos: Fuego 6.5, Tierra 1, Aire 3, Agua 6. Fuego (vitalidad, iniciativa, necesidad de sentido y movimiento hacia la acción) y Agua (profundidad emocional, capacidad de resonancia, sensibilidad abierta): ahí está el peso de tu carta. Es el registro desde el que tu energía se activa con más fluidez y sin esfuerzo consciente. Tierra (arraigo en lo concreto, estructura material, ritmo sostenido y contacto con el cuerpo) está poco representado en tu carta. Ese territorio puede pedirte un esfuerzo consciente, especialmente cuando la presión sube. La mayoría de tus planetas están en signos mutables: hay facilidad para adaptarte, cambiar y responder a lo que ocurre alrededor. Sueles moverte bien cuando las situaciones evolucionan o requieren flexibilidad. La dificultad aparece cuando todo cambia al mismo tiempo y cuesta mantener una dirección estable durante suficiente tiempo. Tu Medio Cielo en Escorpio: La capacidad de trabajar con situaciones complejas, tensas o delicadas sin apartar la mirada. Tu presencia pública puede percibirse como intensa, reservada y difícil de ignorar.

\vspace{0.8cm}
\Needspace{5\baselineskip}

\subsection{2. Ejes Principales}

\subsubsection*{Sol en Aries — Casa 2}
Tu Sol en Aries necesita pasar a la acción antes de que todo esté listo. Tu energía se enciende sola ante un reto o un proyecto nuevo, y puede apagarse si no hay movimiento. Sin algo en marcha o una dirección clara, puedes sentir irritación, vacío o necesidad de activar algo rápidamente. El riesgo no es la falta de energía sino arrancar sin base suficiente y quedarte vacío después del impulso. En Casa 2, necesitas consolidar tu identidad a través del contacto con tus propios recursos, valores y capacidades concretas. Tu seguridad interior suele estar ligada a lo que puedes generar por tus propios medios. Necesitas sentir que puedes sostenerte con recursos propios y confiar en lo que eres capaz de construir de forma autónoma: necesitas una base propia construida desde adentro.

\subsubsection*{Luna en Sagitario — Casa 10}
Tu Luna en Sagitario necesita movimiento, perspectiva y sentido. Cuando algo te afecta, te ayuda entender para qué sirve, hacia dónde te lleva o qué puedes hacer con ello. Si no encuentras una salida, el estado emocional puede convertirse en inquietud, impaciencia o necesidad de cambiar de escenario. El movimiento, el cambio de perspectiva y la sensación de amplitud te ayudan a recuperar estabilidad. En Casa 10, tu estado emocional suele estar conectado con el espacio vocacional y el reconocimiento público. Los logros y los reveses profesionales afectan mucho a cómo te sientes. La visibilidad pública puede ser tanto un apoyo como una fuente de exposición.

\subsubsection*{Ascendente Acuario}
Tu Ascendente en Acuario hace que los demás te perciban como alguien independiente, algo distante y difícil de encajar en lo esperado. Sueles observar primero y participar cuando sientes suficiente espacio o libertad. La primera impresión puede parecer fría o racional, aunque internamente haya mucha más sensibilidad de la que se muestra al principio. Su regente es Urano, en Escorpio y Casa 9. El Ascendente toma el tono de Escorpio: reservada, intensa y difícil de leer rápidamente. Esa energía opera principalmente desde el aprendizaje, la búsqueda de perspectiva y la necesidad de ampliar horizontes.

\subsubsection*{Grados finales de signo}
Los grados finales de signo suelen hablar de zonas de la carta donde hay mucha experiencia acumulada y sensación de       cierre de etapa. La energía no se vive de forma ligera o automática: suele haber más intensidad, más consciencia y la sensación de que una forma antigua de funcionar ya no termina de encajar igual.

No significa algo negativo ni indica un destino fijo. Simplemente muestra lugares donde la vida suele pedir más presencia, más maduración y una manera distinta de responder.

En esta carta aparecen posiciones en grado anarético (29°). Estos grados suelen sentirse como puntos de máxima intensidad dentro de un signo. Hay mucha experiencia acumulada, pero también sensación de presión interna, cansancio de repetir ciertos patrones o necesidad de cambio.

A veces se vive como si ya no fuera posible seguir respondiendo exactamente igual que antes.

Júpiter en Cáncer 29°42', Casa 6 aparece en grado anarético. La necesidad de crecer, expandirte y encontrar sentido aparece con mucha fuerza. Puede costarte permanecer mucho tiempo en experiencias que sientes pequeñas, limitadas o vacías de propósito. A veces puedes buscar siempre el siguiente horizonte antes de terminar de integrar lo que ya estás viviendo. El aprendizaje está en encontrar amplitud sin perder estabilidad.

Medio Cielo en Escorpio 29°12' aparece en grado anarético. La relación con tu dirección profesional y con el lugar que ocupas en el mundo se vive con mucha intensidad. Puede haber una sensación de exigencia interna respecto a lo que vienes a construir o aportar. No suele bastar con funcionar hacia fuera: necesitas sentir que lo que haces tiene sentido real para ti. Cuando esa coherencia falta, el cansancio o la presión pueden sentirse muy profundos.

Fondo del Cielo en Tauro 29°12' aparece en grado anarético. La sensación de hogar, raíz y seguridad interna no suele construirse de forma simple. Puede haber procesos largos de búsqueda de estabilidad emocional o de sensación de pertenencia. El descanso real no depende solo del espacio físico, sino de sentir que existe un lugar interno donde puedas bajar la guardia. Cuando esa base falla, el desorden emocional puede sentirse muy profundo.

También aparecen posiciones en grados previos al anarético. No suelen sentirse tan extremos como el grado 29, pero sí muestran zonas donde ya existe bastante maduración y donde empiezan a aparecer señales de cambio o agotamiento de una forma antigua de funcionar.

Mercurio en Piscis 28°57', Casa 2 aparece en grado previo al anarético. La mente puede funcionar con mucha intensidad y actividad interna. Puede haber necesidad de entender, analizar, nombrar o anticipar lo que ocurre antes de poder relajarte. Pensar puede convertirse en una gran herramienta, pero también en una fuente de sobrecarga cuando intentas resolverlo todo desde la cabeza. Necesitas espacios donde no haga falta entenderlo todo inmediatamente.

Cuando varios puntos de la carta están al final de signo, la vida puede sentirse más intensa o vivirse con menos sensación de margen interno. A veces aparece urgencia, cansancio de ciertas dinámicas o sensación de estar cerca de cambios importantes incluso aunque externamente todo siga igual.

No significa que la vida vaya a ser más difícil. Significa que esas zonas necesitan más consciencia, más pausa y una forma más presente de responder, en lugar de seguir funcionando únicamente desde hábitos automáticos.


\subsubsection*{Signos interceptados}
Los signos Aries (Casa 2) con Sol, Marte y Libra (Casa 8) con Plutón aparecen interceptados: su expresión puede sentirse menos inmediata o menos natural al principio. Esto influye especialmente en Sol, Marte, Plutón, que pueden necesitar más tiempo, experiencia o situaciones concretas para sentirse plenamente desarrollados. Muchas veces estas posiciones no se muestran rápido desde fuera, pero ganan fuerza y claridad con los años.

\vspace{0.8cm}
\Needspace{5\baselineskip}

\subsection{3. Organización de la Energía}

Tu Mercurio en Piscis percibe más por imágenes, sensaciones y asociaciones que por lógica lineal. Puedes captar matices sutiles, pero no siempre resulta fácil explicarlos con orden. Te ayuda dar forma poco a poco a lo que percibes antes de intentar comunicarlo. En Casa 2, la mente tiende a centrarse en recursos, estabilidad y formas concretas de sostenerte. Piensas mucho en lo que tiene valor, en cómo generar seguridad y en qué merece realmente tu energía. Tu Venus en Piscis se vincula con mucha sensibilidad y apertura. Puedes entregarte con facilidad cuando sientes conexión, pero también adaptarte demasiado a la otra persona. El riesgo aparece cuando la compasión, el deseo de unión o la idealización borran tus propios límites. En Casa 1, la forma de relacionarte se percibe rápidamente desde fuera. El vínculo, la atracción y la manera de conectar forman parte importante de cómo te muestras al mundo. Tu Marte en Aries actúa rápido, con iniciativa y poca espera. Cuando algo se activa, necesitas moverte o intervenir. El riesgo aparece cuando la acción se adelanta a la evaluación y terminas gastando fuerza antes de tener una dirección clara. En Casa 2, la acción se dirige a construir estabilidad, recursos y autonomía. Hay impulso real para defender lo propio y sostenerte por tus propios medios.

\vspace{0.8cm}
\Needspace{5\baselineskip}

\subsection{4. Estructura y Tensión}

\subsubsection*{Concentración de energía}

Hay una concentración notable en Casa 2: Sol, Mercurio y Marte. Cuando varios planetas ocupan la misma casa, esa área absorbe buena parte de la actividad vital. En esta carta, el foco de los recursos, la estabilidad y el valor propio funciona como núcleo central: ahí se acumula la presión, la demanda es mayor y también hay más recursos disponibles para operar. No es necesariamente un conflicto, pero sí un punto de alta intensidad sostenida.

También hay una concentración significativa en Casa 7: Saturno, Nodo Norte y Lilith. En este caso, el foco se desplaza hacia las relaciones importantes y el vínculo con otra persona. Es un ámbito donde la exigencia aumenta, donde se concentra parte de la presión y donde también existe capacidad real de sostener procesos importantes. Tampoco es en sí un problema, pero sí un área de intensidad mantenida.



Los aspectos aparecen organizados en dos niveles. Los aspectos exactos (orbe igual o menor a 1°) suelen sentirse de forma más inmediata y constante en la vida cotidiana. Los aspectos estructurales tienen un orbe más amplio, pero siguen describiendo dinámicas importantes que se repiten a lo largo de la vida.



\vspace{0.8cm}
\Needspace{5\baselineskip}
\subsubsection*{Aspectos exactos}
\Needspace{3\baselineskip}
\subsubsection*{Neptuno (Casa 11) trígono Lilith (Casa 7) --- orbe 0.28\textdegree}
El trígono Neptuno-Lilith puede dar acceso a partes muy intuitivas, creativas o difíciles de nombrar. Lo que no encaja fácilmente en la norma puede encontrar salida a través de imágenes, arte, sensibilidad o percepción sutil. El riesgo está en dejarlo demasiado difuso. Te ayuda darle una forma concreta para que no se quede solo como sensación.

\Needspace{3\baselineskip}
\subsubsection*{Marte (Casa 2) quincuncio Saturno (Casa 7) --- orbe 0.53\textdegree}
El quincuncio Marte-Saturno señala un ajuste constante entre impulso y límite. Puedes sentir ganas de actuar y, al mismo tiempo, encontrarte con condiciones que obligan a frenar, ordenar o esperar. No es un bloqueo simple: es una tensión que pide calibrar cuánto empujar y cuándo sostener. Cuando se trabaja bien, permite pasar del impulso rápido a una acción más estable y precisa.

\Needspace{3\baselineskip}
\subsubsection*{Urano (Casa 9) cuadratura Lilith (Casa 7) --- orbe 0.67\textdegree}
La cuadratura Urano-Lilith puede activar reacciones bruscas cuando algo contenido durante demasiado tiempo ya no puede seguir oculto. Puede aparecer como necesidad repentina de romper, cortar, decir algo o salir de una situación. No siempre llega de forma gradual: a veces se manifiesta cuando la tensión ya lleva tiempo acumulándose. Te ayuda detectar antes qué estás tolerando en silencio.

\Needspace{3\baselineskip}
\subsubsection*{Mercurio (Casa 2) trígono Júpiter (Casa 6) --- orbe 0.75\textdegree}
El trígono Mercurio-Júpiter facilita conectar detalles con visión amplia. Puedes comprender el sentido general de una situación y explicarlo de forma que otras personas lo entiendan. Es un buen aspecto para aprender, enseñar, escribir o transmitir ideas. El riesgo aparece cuando ves tan rápido el conjunto que dejas sin ordenar algunos detalles importantes.

\Needspace{3\baselineskip}
\subsubsection*{Plutón (Casa 8) quincuncio Nodo Sur (Casa 1) --- orbe 0.8\textdegree}
El quincuncio Plutón-Nodo Sur señala una tensión entre patrones conocidos y procesos intensos que piden cambio. Puedes permanecer demasiado tiempo en situaciones densas porque, aunque desgasten, resultan familiares. La salida no siempre es romper de golpe; a veces empieza por reconocer cuándo algo ya no transforma, sino que absorbe demasiada energía.

\Needspace{3\baselineskip}
\subsubsection*{Luna (Casa 10) cuadratura Saturno (Casa 7) --- orbe 0.85\textdegree}
La cuadratura Luna-Saturno marca una tensión entre lo que necesitas sentir y lo que te permites mostrar. Puede aparecer como contención emocional, dificultad para pedir apoyo o tendencia a exigirte estar bien antes de recibir cuidado. Desde fuera quizá pareces más fuerte de lo que realmente estás. La clave está en reconocer cuándo la exigencia interna te está costando más que la situación real.

\Needspace{3\baselineskip}
\subsubsection*{Luna (Casa 10) quincuncio Quirón (Casa 3) --- orbe 0.87\textdegree}
El quincuncio Luna-Quirón señala un ajuste delicado entre lo que necesitas emocionalmente y una zona sensible que se activa con facilidad. Puedes sentir que lo que te calma en un momento no sirve en otro, o que ciertas situaciones despiertan una vulnerabilidad difícil de explicar. La clave está en observar qué activa esa sensibilidad antes de intentar resolverla demasiado rápido.

\vspace{0.8cm}
\Needspace{5\baselineskip}
\subsubsection*{Aspectos estructurales}
\Needspace{3\baselineskip}
\subsubsection*{Venus (Casa 1) cuadratura Neptuno (Casa 11) --- orbe 1.15\textdegree}
La cuadratura Venus-Neptuno puede hacer que idealices el vínculo antes de verlo con claridad. Puedes percibir posibilidades, belleza o profundidad donde todavía no hay suficiente realidad concreta. El riesgo aparece cuando sostienes una imagen de la relación más que la relación tal como es. Te ayuda mirar los hechos sin perder sensibilidad.

\Needspace{3\baselineskip}
\subsubsection*{Luna (Casa 10) trígono Marte (Casa 2) --- orbe 1.37\textdegree}
El trígono Luna-Marte facilita que lo que sientes encuentre salida en la acción. Cuando algo te mueve de verdad, puedes responder con rapidez y convertir el estado emocional en movimiento útil. Este aspecto puede ser un recurso de regulación: moverte, actuar o hacer algo concreto antes de quedarte dentro de la emoción.

\Needspace{3\baselineskip}
\subsubsection*{Venus (Casa 1) trígono Urano (Casa 9) --- orbe 1.53\textdegree}
El trígono Venus-Urano da una forma relacional independiente y poco convencional. Sueles necesitar libertad dentro de los vínculos y te atraen relaciones donde haya autenticidad, diferencia o espacio propio. Lo demasiado previsible puede apagarte, pero lo inestable tampoco siempre te sostiene. La clave está en diferenciar libertad real de distancia emocional.

\Needspace{3\baselineskip}
\subsubsection*{Saturno (Casa 7) trígono Quirón (Casa 3) --- orbe 1.71\textdegree}
El trígono Saturno-Quirón indica que la estructura puede ayudarte a sostener zonas sensibles de tu vida. Cuando tienes ritmo, límites y responsabilidades claras, manejas mejor aquello que normalmente te hace sentir inseguridad o exposición. No se trata de endurecerte, sino de crear un contenedor suficientemente estable para no quedarte en la vulnerabilidad.

\Needspace{3\baselineskip}
\subsubsection*{Sol (Casa 2) cuadratura Júpiter (Casa 6) --- orbe 4.19\textdegree}
La cuadratura Sol-Júpiter puede llevarte a ampliar más de lo que puedes sostener. Hay entusiasmo, visión y capacidad de crecimiento, pero también tendencia a asumir demasiado o confiar en que podrás con todo. El límite no viene a apagar la expansión, sino a darle una forma que no te desgaste.

\Needspace{3\baselineskip}
\subsubsection*{Sol (Casa 2) trígono Neptuno (Casa 11) --- orbe 5.17\textdegree}
El trígono Sol-Neptuno aporta sensibilidad, imaginación y capacidad para percibir matices que no son evidentes. Puede favorecer lo creativo, lo contemplativo o el acompañamiento de otras personas. El riesgo aparece si esa apertura no se traduce en algo concreto. Te ayuda dar forma práctica a lo que percibes.

\Needspace{3\baselineskip}
\subsubsection*{Sol (Casa 2) oposición Plutón (Casa 8) --- orbe 7.88\textdegree}
La oposición Sol-Plutón marca una tensión intensa entre afirmarte y atravesar procesos que no puedes controlar del todo. Puede haber momentos en los que la vida te obliga a revisar quién eres, qué deseas y qué poder estás entregando o reteniendo. No es una tensión ligera: suele pedir honestidad, límites y capacidad de sostener cambios importantes sin destruirte por dentro.



\vspace{0.8cm}
\Needspace{5\baselineskip}

\subsection{5. Sostén y Equilibrio}

Tu Quirón en Tauro, Casa 3, señala una zona sensible relacionada con la estabilidad, el valor propio y la relación con lo que necesitas para sostenerte. Puede que hayas vivido inseguridad material o sensación de no merecer descanso, placer o apoyo suficiente. Esto suele notarse especialmente en la comunicación, el aprendizaje y el entorno cercano. No significa que tengas límites en ese terreno, pero sí que probablemente necesites más tiempo, cuidado o consciencia para sentir estabilidad ahí.

\vspace{0.3cm}
Tu Nodo Sur en Piscis, Casa 1, muestra una tendencia a moverte desde la falta de límites claros, la adaptación excesiva o la tendencia a perderte en el entorno, especialmente en tu forma de afirmarte y mostrarte. Suele ser una forma de funcionar conocida o automática para ti. Tu Nodo Norte en Virgo, Casa 7, señala la necesidad de desarrollar orden, presencia en lo cotidiano y capacidad de concretar, especialmente en las relaciones importantes y el vínculo con otra persona. No se trata de rechazar lo anterior, sino de ampliar tu forma habitual de responder.

\vspace{0.3cm}
En tu carta hay poca presencia de Tierra. Cuando aparece saturación, bloqueo o sensación de desorganización interna, suele ayudarte volver a espacios relacionados con el contacto con el cuerpo, la rutina y lo concreto. Eso puede darte más estabilidad y ayudarte a recuperar claridad.

\vspace{0.8cm}
\Needspace{5\baselineskip}

\subsection{6. Orientación}

\subsubsection*{Qué observar}
Tu Ascendente en Acuario describe la forma en que otras personas suelen percibirte al principio. Eso no siempre coincide con cómo vives las cosas internamente o con las tensiones más importantes de tu carta. Cuando hay demasiada distancia entre lo que muestras y lo que realmente te ocurre, aparece desgaste.

\subsubsection*{Qué puede desestabilizar}
Hay ciertas tensiones en tu carta que pueden activarse con bastante fuerza: Nodo Norte-Nodo Sur, Marte-Saturno, Urano-Lilith. Cuando varias coinciden al mismo tiempo, puede aparecer sensación de saturación, reacción excesiva o dificultad para mantener claridad.

\subsubsection*{Qué estructura y sostiene}
Tu Sol en Aries, Casa 2, muestra un terreno importante de estabilidad para ti. Cuando tu forma de actuar está alineada con lo que realmente valoras, suele aparecer más claridad y menos dispersión.

\subsubsection*{Orientación práctica}
Una orientación práctica: en tu carta hay poca Tierra. Cuando aparece saturación o exceso de actividad mental, suele ayudarte volver a cosas muy concretas y físicas: mover el cuerpo, cocinar, ordenar, caminar o sostener una rutina sencilla durante varios días seguidos. Lo importante no es hacerlo perfecto, sino recuperar sensación de estabilidad.

\vspace{1cm}
\begin{center}
{\small\itshape\color{grisai}
La astrología se usa aquí como lenguaje simbólico de observación, no como una definición de quién eres.\\
Esta interpretación es tu mapa de posibilidades, no tu destino fijo.
}
\end{center}

\end{document}
